\section{Discussion}
\paragraph{Connection of the Taxonomy with SDLC Stages}
The proposed interaction taxonomy aligns closely with different stages of the Software Development Lifecycle (SDLC), showing that developer-AI collaboration varies according to task goals and abstraction levels. Conversational assistance is most relevant during requirements elicitation and ideation, where developers use natural language interaction to clarify ambiguity and explore solutions. Suggestion-based and generative assistance primarily support implementation tasks by accelerating coding through completions and automated code generation. Retrieval and explanation assistance are more closely associated with debugging and review activities, where LLMs help developers interpret errors, summarize documentation, and reason about program behavior. Documentation workflows often combine multiple interaction modes, demonstrating that AI-assisted programming is fundamentally task-dependent across SDLC stages.

\paragraph{RQ 1: How do Generative AI tools reshape developers' problem solving strategies during software development?}
Generative AI shifts developer problem-solving from manual solution construction toward solution exploration, evaluation, and refinement. Across the surveyed studies, developers increasingly interact with AI systems through conversational prompts, code completions, and iterative refinement workflows rather than writing all code from scratch. Suggestion-based tools accelerate familiar tasks through autocomplete-like interactions, while conversational and generative systems support exploratory reasoning in unfamiliar domains. This changes the developer's role from primarily authoring code to reviewing, validating, and adapting AI-generated outputs. Studies such as \citet{barke2022grounded} and \citet{bird2023copilot} further show that developers dynamically switch between \textit{"acceleration"} and \textit{"exploration"} modes depending on task familiarity, with exploration involving deeper cognitive engagement through prompt iteration, comparison of alternatives, and reasoning about generated solutions.

\paragraph{RQ 2: How does AI-assisted development influence developer behavior across different stages of the SDLC?}
AI-assisted development influences developer strategies differently across SDLC stages by aligning distinct interaction modes with specific development activities. During requirements elicitation, conversational AI supports exploratory reasoning, clarification of ambiguous requirements, and structured artifact generation through natural language interaction. In coding, suggestion based and generative assistance accelerate implementation by automating syntax completion, boilerplate generation, and \textit{natural language to code} translation. During debugging and review, retrieval and explanation assistance help developers understand errors, summarize documentation, and reason about program behavior, although current systems still struggle with complex logical and multi-file defects. Documentation workflows benefit from conversational and generative capabilities that automatically produce explanations, comments, and summaries.


\paragraph{RQ 3: What productivity benefits and tradeoffs are associated with AI-assisted software development?}
AI-assisted development provides substantial productivity gains, particularly by reducing implementation effort and accelerating task completion. Empirical evidence from \citet{peng2023copilot} shows that developers using GitHub Copilot completed programming tasks 55.8\% faster than those without AI assistance, while other studies report improvements in code generation, documentation creation, and debugging support. AI systems also reduce cognitive barriers for unfamiliar technologies by enabling iterative learning and exploratory interaction.

However, these benefits come with important costs. Developers often spend more time reviewing and validating generated code, which can reduce deep code understanding and increase cognitive load. Overreliance on AI suggestions may lead to anchoring bias, misplaced trust, and acceptance of incorrect outputs. Additionally, current LLMs struggle with logical reasoning, edge cases, large multi-file contexts, and complex debugging tasks, making human oversight essential. Thus, while AI-assisted development significantly improves efficiency and accessibility, it introduces tradeoffs in reliability, comprehension, and developer control.